% concatenated from Rigidity_and_volume_pinching_for_the_sharp_gradient_estimate_in_positive_Ricci_curvature.tex
\documentclass[11pt]{amsart}
\usepackage{amsmath,amssymb,amsthm}
\usepackage{mathtools}
\usepackage{esint}
\usepackage{comment}
\allowdisplaybreaks

\usepackage{xcolor}
\usepackage{hyperref}

\theoremstyle{plain}
\newtheorem{theorem}{Theorem}[section]
\newtheorem{remark}[theorem]{Remark}
\newtheorem{definition}[theorem]{Definition}
\newtheorem{corollary}[theorem]{Corollary}
\newtheorem{lemma}[theorem]{Lemma}
\newtheorem{proposition}[theorem]{Proposition}

\numberwithin{equation}{section}

\newcommand{\R}{\mathbb{R}}
\newcommand{\C}{\mathbb{C}}
\newcommand{\N}{\mathbb{N}}
\newcommand{\Z}{\mathbb{Z}}
\newcommand{\Q}{\mathbb{Q}}
\renewcommand{\S}{\mathbb{S}}
\newcommand{\RP}{\mathbb{RP}}

\newcommand{\ve}{\varepsilon}
\newcommand{\E}{\mathcal{E}}
\newcommand{\wh}[1]{\widehat{#1}}
\newcommand{\norm}[1]{\left\|#1\right\|}
\newcommand{\hess}{\mathrm{Hess}}
\DeclareMathOperator{\ric}{Ric}
\DeclareMathOperator{\vol}{Vol}
\DeclareMathOperator*{\argmax}{arg\,max}

%blow-up quantity
\newcommand{\hd}{\widehat{\Delta}}
\newcommand{\hn}{\widehat{\nabla}}
\newcommand{\hg}{\widehat{g}}
\newcommand{\hw}{\widehat{W}}
\newcommand{\hm}{\widehat{M}}
\newcommand{\dv}{\,dV_{\widehat{g}}}
\newcommand{\da}{\,dA_{\widehat{g}}}

\newcommand{\dvol}{\,dV_g}
\newcommand{\darea}{\,dA_g}

%sublevel set
\newcommand{\sls}{\{\rho\le r\}}
\newcommand{\ls}{\{\rho=r\}}
\newcommand{\bsls}{\{b\le r\}}
\newcommand{\bls}{\{b=r\}}

\title[Rigidity and volume pinching]{Rigidity and volume pinching for the sharp gradient estimate in positive Ricci curvature}

\author[Kim]{Junyoung Kim}
\address{Korea Advanced Institute of Science and Technology (KAIST)\\
Department of Mathematical Sciences\\
Daejeon, South Korea}
\email{junykim@kaist.ac.kr}

\author[Park]{Jiewon Park}
\address{Korea Advanced Institute of Science and Technology (KAIST)\\
Department of Mathematical Sciences\\
Daejeon, South Korea}
\email{jiewonpark@kaist.ac.kr}

\begin{document}

\begin{abstract}
Colding established a sharp gradient estimate for the Green function on manifolds with nonnegative Ricci curvature, which was extended to positive Ricci curvature by Manea recently. We prove almost rigidity of such gradient estimates for positive Ricci curvature. We show that the average of the gradient deficit, namely $\fint (1-|\nabla b|^2)\,dV$, controls the volume deficit in a quantitative manner; here $b$ is the distance-like function defined using the Green function. This implies a quantitative almost rigidity theorem for manifolds with $\ric\ge(n-1)g$ and a gap theorem for Einstein manifolds. We also prove a rigidity theorem for closed 4-dimensional Einstein manifolds by finding a new monotonicity formula.
\end{abstract}

\maketitle

\section{Introduction}

Gradient estimates for Green functions provide a useful bridge between analysis and geometry. On complete Riemannian manifolds with nonnegative Ricci curvature, Colding introduced a distance-like function associated with the Green function and proved a sharp gradient estimate together with monotonicity formulas measuring the failure of the manifold to be conical \cite{Col12}. Recently, Manea obtained the positive-Ricci counterpart on closed manifolds \cite{Man25}.

In this paper we are also concerned with the positive-Ricci setting. Let $(M^n,g)$ be a closed connected Riemannian manifold of dimension $n\ge 3$ with $\ric\ge(n-1)g$. For $p\in M$, let $G_p$ be the positive Green function of the operator
\begin{equation*}
L=-\Delta+\frac{n(n-2)}{4},
\end{equation*}
with the normalization
\begin{equation}\label{eq:green-function-normalization}
LG_p=(n-2)\omega_{n-1}\delta_p,
\end{equation}
where $\omega_{n-1}=\vol(\S^{n-1})$. The function
\begin{equation}\label{eq:function b}
b_p=2\arcsin\left(\frac{1}{2}G_p^{\frac{1}{2-n}}\right),\quad b_p(p)=0,
\end{equation}
is well-defined, takes values in $[0,\pi]$, and satisfies the sharp estimate
\begin{equation}\label{eq:manea}
    |\nabla b_p|\le 1\quad\text{on $M\setminus\{p\}$}
\end{equation}
by \cite{Man25}. On the round sphere, $b_p$ is exactly the distance from $p$, so equality holds throughout the regular set. Conversely, equality at one point in \eqref{eq:manea} characterizes the round sphere \cite{Man25}.

We ask the natural question whether a suitable lower bound on $|\nabla b_p|$ would imply closeness to the round sphere. Obviously, one cannot impose a global positive lower bound on $|\nabla b_p|$, simply because the manifold is compact and $|\nabla b_p(x)| \to 1$ as $x \to p$ (cf. \cite[Lemma 2.4]{Man25}). Hence we consider the normalized integral gap
\begin{equation}\label{eq:integral gap}
\E_p(M,g):=\fint_M(1-|\nabla b_p|^2)\,dV_g,
\end{equation}
which, by \eqref{eq:manea}, always satisfies $0\le\E_p(M,g)\le 1$. The definition \eqref{eq:integral gap} permits the defect $(1-|\nabla b_p|^2)$ to concentrate on small regions. It is therefore not immediate that smallness of $\E_p(M,g)$ should control any global geometric quantity.

In this paper we prove two main theorems, obtained by two different methods. The first is a rigidity theorem for the Weyl tensor of closed Einstein 4-manifolds under a lower bound on the gradient, proved through a new monotonicity formula for the ratio of two Green functions. The second is a quantitative volume pinching theorem, valid in every dimension $n\ge 3$ under a Ricci lower bound. It shows that the normalized integral gap controls the volume deficit. As consequences of the second theorem, we deduce an almost rigidity theorem for manifolds with $\ric\ge(n-1)g$ and a gap theorem for Einstein manifolds.

For our first main result, we consider $(M^4,g)$ a closed connected Einstein 4-manifold with $\ric=3g$. Fix distinct points $p,q\in M$, and consider the Green-function blow-up 
\begin{equation}\label{eq:blow-up-data}
\hg=G_q^2g\quad\text{on $\hm=M\setminus\{q\}$}\quad\text{and}\quad \rho=\left(\frac{G_q}{G_p}\right)^{1/2}.
\end{equation}
The metric $\hg$ is complete, scalar-flat, and asymptotically flat, while $\rho$ runs from $0$ at $p$ to infinity along the end corresponding to $q$. Thus $\rho$ plays the role of a radial function on the blow-up.

\begin{theorem}[Four-dimensional Weyl rigidity]\label{thm1}
There exists a universal constant $c>0$ with the following property. In the setting of \eqref{eq:blow-up-data}, suppose that, for some $\delta>0$, 
\begin{equation*}
|\hn\rho|_{\hg}\ge\delta\quad\text{on $M\setminus\{p,q\}$}\quad \text{and}\quad \sup_{\hm}\rho^2|\hw|_{\hg}\le c\delta^2,
\end{equation*}
then $\hw\equiv 0$ on $\hm$, and $(M,g)$ is isometric to the round sphere $(\S^4,g_{\S^4})$.
\end{theorem}

The proof of Theorem \ref{thm1} adapts the ODE method of Lee--Park for complete Ricci-flat manifolds \cite{LP25} to a conformal blow-up of a positive Einstein manifold. The new tool is a weighted differential inequality for the density $G_q^{-1}|\hw|^2_{\hg}$, derived from Wu's conformally Einstein Weitzenb\"ock formula \cite{Wu2017}. Integrating this over sublevel sets of $\rho$ produces a function $F(r)$ satisfying an Euler-type differential inequality. The pole expansions at $p$ and at the asymptotically flat end impose incompatible boundary behavior unless $F$ vanishes identically. Consequently, the Weyl tensor vanishes, and the lower bound for $|\hn\rho|_{\hg}$ rules out the nontrivial spherical space form.

Our second main result, Theorem \ref{thm2}, states that the integral gap $\E_p(M,g)$ quantitatively controls the volume deficit. In fact the result holds in any dimension $n\ge 3$ and requires only a positive Ricci lower bound, not necessarily Einstein.

\begin{theorem}[Volume pinching by the integral gap]\label{thm2}
Let $n\ge 3$ and $0<\gamma<\frac{2}{n+2}$. Then there exists a constant $C_{n,\gamma}>0$ with the following property. If $(M^n,g)$ is a closed connected Riemannian $n$-manifold with $\ric\ge(n-1)g$ and $p\in M$, then 
\begin{equation}\label{eq:thm2}
    1-\frac{\vol(M,g)}{\vol(\S^n)}\le C_{n,\gamma}\E_p(M,g)^\gamma.
\end{equation}
\end{theorem}

We remark that Theorem \ref{thm2} holds just with a lower bound on Ricci curvature; no Einstein assumption and no a priori smallness condition for $\E_p(M,g)$ are required in Theorem \ref{thm2}. 
The proof proceeds through a weighted quantity adapted to the Green function. In dimension four, the estimate \eqref{eq:thm2} holds with the exponent $\gamma<\frac{1}{2}$; see Remark \ref{rmk3.5}.

Theorem \ref{thm2} has two corollaries.  In the setting of nonnegative Ricci curvature, Honda--Peng \cite{HondaPeng} established the almost rigidity of the gradient estimate. They showed that under the pointwise assumption $|\nabla b|\ge 1-\delta$, $M$ is measured Gromov-Hausdorff close to the Euclidean space; in fact their results hold more generally for non-parabolic $\mathsf{RCD}(0,N)$ spaces. Our first corollary is an analogous almost rigidity theorem in the positive-Ricci setting, in which only the smallness of the normalized integral gap \eqref{eq:integral gap} is assumed.

\begin{corollary}[Almost rigidity for Ricci lower bound]\label{cor3}
For $n\ge 3$, there exist constants $\ve=\ve(n)>0,C_n>0$ and $\alpha_n>0$ with the following property. Let $(M^n,g)$ be a closed connected Riemannian $n$-manifold with $\ric\ge(n-1)g$ and $p\in M$. If 
\begin{equation*}
\E_p(M,g)\le \ve,
\end{equation*}
then $M$ is diffeomorphic to the round sphere $(\S^n,g_{\S^n})$ and
\begin{equation*}
d_{\mathrm{GH}}((M,g),(\S^n,g_{\S^n}))\le C_n\E_p(M,g)^{\alpha_n}.
\end{equation*}
\end{corollary}

Indeed, once we obtain the almost-maximal volume of Theorem \ref{thm2}, the quantitative volume pinching theorem of Aubry \cite[Theorem 4]{Aubry2005} yields both the diffeomorphism and a power-type Gromov--Hausdorff estimate in terms of the volume deficit. The diffeomorphism can also be obtained using Cheeger--Colding theory \cite{CheegerColding1997}. One may take $\alpha_n=\gamma\beta(n)$ for any $\gamma<\frac{2}{n+2}$ and $\beta(n)$ the exponent of \cite{Aubry2005}.


For Einstein metrics, almost-maximal volume implies isometry to the round sphere. Combining Theorem \ref{thm2} with the sphere theorem for Einstein manifolds of Honda--Mondello \cite[Theorem C]{HondaMondello}, we obtain the following result; see also \cite[Theorem 3.12]{HondaMondino}, and \cite[Chapter 12]{Besse} for an alternative proof via the moduli space of Einstein metrics.

\begin{corollary}\label{thm3}
For $n\ge 3$, there exists a constant $\ve=\ve(n)>0$ with the following property. Let $(M^n,g)$ be a closed connected Einstein $n$-manifold with $\ric=(n-1)g$ and $p\in M$. If
\begin{equation*}
\E_p(M,g)\le \ve,
\end{equation*}
then $(M,g)$ is isometric to the round sphere $(\S^n,g_{\S^n})$.
\end{corollary}

Finally, the hypotheses of Theorem \ref{thm1} are automatically satisfied when the normalized integral gap is sufficiently small, provided the second pole $q$ is chosen to be the maximum point of $b_p$.

\begin{theorem}\label{thm4}
Let $0<\delta<2$ and $\zeta>0$. Then there exists a constant $\ve_0=\ve_0(\delta,\zeta)>0$ with the following property. Let $(M^4,g)$ be a closed connected Einstein manifold with $\ric=3g$, fix $p\in M$, and choose 
\begin{equation*}
q\in\argmax_{x\in M}b_p(x).
\end{equation*}
If
\begin{equation*}
\E_p(M,g)\le\ve_0,
\end{equation*}
then
\begin{equation*}
|\hn\rho|_{\hg}\ge\delta\quad\text{on $M\setminus\{p,q\}$}\quad\text{and}\quad\sup_{\hm}\rho^2|\wh{W}|_{\hg}<\zeta.
\end{equation*}
\end{theorem}

Taking $\zeta<c\delta^2$ in Theorem \ref{thm4} and applying Theorem \ref{thm1} gives a second proof of Corollary \ref{thm3} in dimension four in the case $q\in\argmax_{x\in M}b_p(x)$, which does not rely on the sphere theorem of \cite{HondaMondello}.

\subsection*{Organization of the paper}
In Section 2, we prove the four-dimensional Weyl rigidity theorem by deriving the weighted Weitzenb\"ock inequality, establishing the endpoint asymptotics of the sublevel-set functional, and applying an ODE argument by finding a monotone functional. In Section 3, we develop the level-set identities, prove the quantitative volume pinching, and deduce the Einstein gap theorem. In Section 4, we verify the blow-up hypotheses under small integral gap and obtain another four-dimensional proof without using the sphere theorem.

\subsection*{Acknowledgment} This work was supported by the National Research Foundation of Korea(NRF) grant funded by the Korea government(MSIT) RS-2024-00346651. 

\section{Proof of Theorem \ref{thm1}}
Throughout this section, $(M^4,g)$ is a closed connected Einstein 4-manifold with $\ric_g=3g$. For $a\in M$, we denote by $G_a$ the positive Green function for $L=-\Delta+2$ with pole at $a$ and normalization $LG_a=2\omega_3\delta_a$. Fix two distinct points $p,q\in M$ and set
\begin{equation*}
\hg=G_q^2g\quad\text{on $\hm:=M\setminus\{q\}$},\quad \rho=\left(\frac{G_q}{G_p}\right)^{1/2}\quad\text{on $M\setminus\{p,q\}$.}
\end{equation*}
By conformal covariance, $\hg$ is scalar-flat; the standard Green function expansion also shows that it is complete and asymptotically flat \cite{LeeParker1987}. Note that $\rho$ is continuous and positive on $M\setminus\{p,q\}$. Moreover, $\rho$ extends continuously to $p$ with $\rho(p)=0$ and tends to infinity at the asymptotically flat end corresponding to $q$. 

\begin{lemma}\label{lem:rho equations}
On $M\setminus\{p,q\}$, we have
\begin{equation}\label{eq:lem2.1}
\hd\rho=3\rho^{-1}|\hn\rho|^2_{\hg},\qquad \hd\rho^2=8|\hn\rho|_{\hg}^2.
\end{equation}
\end{lemma}

\begin{proof}
The conformal covariance of the conformal Laplacian and the scalar-flatness of $\hg$ give
\begin{equation*}
0=L_{\hg}\left(\frac{G_p}{G_q}\right)=-\hd\rho^{-2}.
\end{equation*}
Since 
\begin{equation*}
\hd\rho^{-2}=-2\rho^{-3}\hd\rho+6\rho^{-4}|\hn\rho|_{\hg}^2,
\end{equation*}
the first identity follows. Then,
\begin{equation*}
\hd\rho^2=2\rho\hd\rho+2|\hn\rho|_{\hg}^2=8|\hn\rho|_{\hg}^2,
\end{equation*}
which proves the second assertion.
\end{proof}

The following lemma replaces the Kato inequality used in \cite{LP25}. Its proof uses the conformally Einstein Weitzenb\"ock formula in \cite{Wu2017}.

\begin{lemma}\label{lem:Kato}
Define $ V:=G_q^{-1}|\hw|^2_{\hg}$. Then there exists a universal constant $C>0$ such that
\begin{equation}\label{eq:lem2.2}
\hd V\ge-C|\hw|_{\hg}V.
\end{equation} 
\end{lemma}

\begin{proof}
Let
\begin{equation*}
f=2\log G_q\quad\text{on $\hm$}
\end{equation*}
so that $g=e^{-f}\hg$. We note that $\ric_g=3g$ implies
\begin{equation*}
\ric_{\hg}+\wh{\hess}f+\frac{1}{2}df\otimes df=\lambda\hg,
\end{equation*}
where
\begin{equation}\label{eq:lem2.2-1}
4\lambda=\hd f+\frac{1}{2}|\hn f|_{\hg}^2.
\end{equation}
On the other hand, the scalar curvature transformation for $g=e^{-f}\hg$ implies 
\begin{equation}\label{eq:lem2.2-2}
\hd f-\frac{1}{2}|\hn f|_{\hg}^2=4G_q^{-2}.
\end{equation}
Combining \eqref{eq:lem2.2-1} and \eqref{eq:lem2.2-2}, we obtain
\begin{equation}\label{eq:lem2.2-3}
4\lambda-|\hn f|_{\hg}^2=4G_q^{-2}.
\end{equation}

Let 
\begin{equation*}
\hd_f=\hd-\langle\hn f,\hn\cdot\rangle.
\end{equation*}
The conformally Einstein Weitzenb\"ock formula in \cite[Corollary 1.1]{Wu2017} states that
\begin{equation*}
\hd_f|\hw^{\pm}|^2_{\hg}=2|\hn\hw^{\pm}|^2_{\hg}+\left(4\lambda-|\hn f|^2_{\hg}\right)|\hw^{\pm}|^2_{\hg}-144\det \hw^{\pm}
\end{equation*}
on an oriented coordinate neighborhood. The estimate $|\det\hw^{\pm}|\le C|\hw^{\pm}|^3_{\hg}$, together with \eqref{eq:lem2.2-3}, yields 
\begin{equation}\label{eq:lem2.2-4}
\hd_f|\hw|^2_{\hg}\ge 4G_q^{-2}|\hw|^2_{\hg}-C|\hw|^3_{\hg}.
\end{equation}
Note that the inequality is independent of the chosen local orientation and is global.

A direct computation gives
\begin{align*}
e^{-f/2}\hd_f|\hw|^2_{\hg}&=e^{-f/2}\hd_f(e^{f/2}V)\\&=\hd V+\left(\frac{1}{2}\hd f-\frac{1}{4}|\hn f|^2_{\hg}\right)V\\&=\hd V+2G_q^{-2}V,
\end{align*}
where we used \eqref{eq:lem2.2-2}. Multiplying \eqref{eq:lem2.2-4} by $e^{-f/2}$ then gives
\begin{equation*}
\hd V+2G_q^{-2}V\ge 4G_q^{-2}V-C|\hw|_{\hg}V,
\end{equation*}
which proves \eqref{eq:lem2.2}.
\end{proof}


We assume that, for some $\delta>0$, 
\begin{equation}\label{eq:assumption sec1}
|\hn\rho|_{\hg}\ge\delta\quad\text{on $M\setminus\{p,q\}$.}
\end{equation}
Every positive value of $\rho$ is then regular. We define the function $F:[0,\infty)\rightarrow\R$ by
\begin{align}
F(r)=\int_{\sls}V|\hn\rho|_{\hg}^2\dv.
\end{align}
Each sublevel set $\{\rho\le r\}$ is a compact subset of $\hm$, so $F$ is well-defined.

We use the following notation for asymptotic expansions. For a function $f(r,\theta)$ in geodesic polar coordinates, we denote by $f=O_k(r^\mu)$ if
\begin{equation*}
|\partial_r^j\nabla_\theta^\alpha f(r,\theta)|\le C_{j,\alpha} r^{\mu-j}\quad (j+|\alpha|\le k),
\end{equation*}
uniformly in $\theta$. The Green function expansion \cite{LeeParker1987} gives
\begin{equation}\label{eq:asymptotic of G}
G_p(r,\theta)=r^{-2}(1+O_1(r^2)).
\end{equation}
By symmetry of the Green function, set
\begin{equation*}
\gamma:=G_q(p)=G_p(q)>0.
\end{equation*}

\begin{lemma}\label{lem:graphical expression}
Let $s=d_g(p,\cdot)$ and $t=d_g(q,\cdot)$. There exist two smooth functions $A_p, A_q\in C^\infty(\S^3)$ such that 
\begin{align}
\label{eq:rho near p} \rho(s,\theta)=\gamma^{1/2}s+A_p(\theta)s^2+O_1(s^3)&\quad \text{as } s\rightarrow0,\\
\label{eq:rho near q} \rho(t,\theta)=\gamma^{-1/2}t^{-1}+A_q(\theta)+O_1(t)&\quad\text{as } t\rightarrow0.
\end{align}
In particular, for all sufficiently small $r>0$, the level set $\ls$ near $p$ can be expressed as the polar graph
\begin{equation*}
s=\sigma(r,\theta),\quad \sigma(r,\theta)=\gamma^{-1/2}r+O_1(r^2),\quad \partial_r\sigma=\gamma^{-1/2}+O(r),
\end{equation*}
and for all sufficiently large $r>0$, the same level set near $q$ can be expressed as the polar graph
\begin{equation*}
t=\tau(r,\theta),\quad \tau(r,\theta)=\gamma^{-1/2}r^{-1}+O_1(r^{-2}),\quad \partial_r\tau=-\gamma^{-1/2}r^{-2}+O(r^{-3}).
\end{equation*}
Moreover, $\sls=\{s\le\sigma(r,\theta)\}$ for all sufficiently small $r$, whereas $\{\rho>r\}=\{0<t<\tau(r,\theta)\}$ for all sufficiently large $r$.
\end{lemma}

\begin{proof}
Noting that $G_q$ is positive and smooth near $p$, the Taylor expansion in the normal coordinates, together with \eqref{eq:asymptotic of G}, gives
\begin{equation*}
G_q(s,\theta)=\gamma+B_p(\theta)s+O_2(s^2),\quad G_p^{-1/2}(s,\theta)=s(1+O_1(s^2))
\end{equation*}
for some constant $\gamma>0$ and some smooth function $B_p\in C^\infty(\S^3)$. Expanding $G_q^{1/2}$ and multiplying $G_p^{-1/2}$ proves \eqref{eq:rho near p}. We can obtain \eqref{eq:rho near q} by using the same argument.

Equation \eqref{eq:rho near p} implies $\partial_s\rho=\gamma^{1/2}+O(s)>0$ for all sufficiently small $s$. Then, the implicit function theorem yields the graph expression. Also, if we substitute the graph $s=\sigma(r,\theta)$ in \eqref{eq:rho near p} and differentiate in $r$, then we have the desired expansions of $\sigma$ and $\partial_r\sigma$. Similarly, equation \eqref{eq:rho near q} gives $\partial_t\rho^{-1}=\gamma^{1/2}+O(t)>0$ for small $t$, and applying the implicit function theorem to $\rho^{-1}$ gives the graph expression $t=\tau(r,\theta)$ and the desired expansions. 
\end{proof}

\begin{lemma}
With the notation of Lemma \ref{lem:graphical expression},
\begin{align}
\label{eq:volume element near p}
V|\hn\rho|_{\hg}^2\dv=\left(\gamma^{-2}|W(p)|^2_gs^3+O(s^4)\right)\,ds\,d\theta&\quad\text{as $s\rightarrow0$},\\\label{eq:volume element near q}
V|\hn\rho|_{\hg}^2\dv=\left(\gamma^{-1}|W(q)|^2_g t^5+O(t^6)\right)\,dt\,d\theta&\quad\text{as $t\rightarrow0$},
\end{align}
where $d\theta$ is the standard volume measure on $\S^3$.
\end{lemma}

\begin{proof}
Under $\hg=G_q^2g$ in dimension four, we have
\begin{equation*}
|\wh{W}|_{\hg}^2=G_q^{-4}|W|_g^2,\quad |\hn \rho|_{\hg}^2=G_q^{-2}|\nabla \rho|_g^2,\quad \dv=G_q^4\,dV_g,
\end{equation*}
and thus
\begin{align}\label{eq:lem2.4-1}
V|\hn\rho|_{\hg}^2\dv=G_q^{-3}|W|^2_g|\nabla\rho|_g^2\,dV_g.
\end{align}

In four-dimensional geodesic polar coordinates centered at either pole $p$ or $q$, we have
\begin{equation}\label{eq:asymptotic of vol}
dV_g=r^{3}(1+O_2(r^2))\,dr\,d\theta.
\end{equation}
Then, at $p$, equations \eqref{eq:rho near p} and \eqref{eq:asymptotic of vol} imply
\begin{equation*}
|\nabla\rho|_g^2=\gamma+O(s),\quad G_q^{-3}=\gamma^{-3}+O(s),\quad |W|^2_g=|W(p)|^2_g+O(s),
\end{equation*}
and substitution in \eqref{eq:lem2.4-1} proves \eqref{eq:volume element near p}. Similarly, we can prove \eqref{eq:volume element near q}.
\end{proof}

\begin{proposition}[Asymptotics of $F$ near $0$]\label{prop:asymptotic of F near 0}
As $r\rightarrow0$, we have
\begin{align}
    \label{eq:asymptotic of F near 0}F(r)&=\frac{\omega_3}{4}\gamma^{-4}|W(p)|^2_g\,r^4+O(r^5),\\
    \label{eq:asymptotic of F' near 0}F'(r)&=\omega_3\gamma^{-4}|W(p)|_g^2\,r^3+O(r^4).
\end{align}
In particular, $F(r)=O(r^4)$ and $F'(r)=O(r^3)$ as $r\rightarrow0$.
\end{proposition}

\begin{proof}
For small $r$, Lemma \ref{lem:graphical expression} and \eqref{eq:volume element near p} give
\begin{equation*}
F(r)=\int_{\S^3}\int_0^{\sigma(r,\theta)}\left(\gamma^{-2}|W(p)|^2_g\,s^3+O(s^4)\right)\,ds\,d\theta,
\end{equation*}
where the remainder is uniform in $\theta$. Since $\sigma(r,\theta)=\gamma^{-1/2}r+O(r^2)$, we have
\begin{align*}
F(r)&=\int_{\S^3}\left(\frac{\gamma^{-2}|W(p)|_g^2}{4}\sigma(r,\theta)^4+O(\sigma(r,\theta)^5)\right)\,d\theta\\
&=\frac{\omega_3}{4}\gamma^{-2}|W(p)|^2_g\left(\gamma^{-2}r^4+O(r^5)\right)+O(r^5),
\end{align*}
which is \eqref{eq:asymptotic of F near 0}. 

On the other hand, differentiating $F$ in $r$ yields that
\begin{align*}
F'(r)&=\int_{\S^3}\left(\gamma^{-2}|W(p)|^2_g\sigma^3+O(\sigma^4)\right)\partial_r\sigma \,d\theta\\
&=\omega_3\gamma^{-2}|W(p)|^2_g\left(\gamma^{-3/2}r^3+O(r^4)\right)\left(\gamma^{-1/2}+O(r)\right)+O(r^4)\\
&=\omega_3\gamma^{-4}|W(p)|_g^2r^3+O(r^4),
\end{align*}
which proves \eqref{eq:asymptotic of F' near 0}.
\end{proof}

\begin{proposition}[Asymptotics of $F$ as $r\rightarrow\infty$]\label{prop:asymptotic of F at infinity}
The integral
\begin{equation*}
F_\infty:=\int_{\hm}V|\hn\rho|_{\hg}^2\dv
\end{equation*}
is finite, and $F(r)\rightarrow F_\infty$ as $r\rightarrow\infty$. More precisely,
\begin{equation}
\label{eq:asymptotic of F at infinity}
F_\infty-F(r)=\frac{\omega_3}{6}\gamma^{-4}|W(q)|_g^2\,r^{-6}+O(r^{-7})\quad\text{as } r\rightarrow\infty.
\end{equation}
In particular, $F_\infty-F(r)=O(r^{-6})$ as $r\rightarrow\infty$. 
\end{proposition}

\begin{proof}
The volume element has the forms $O(s^3)ds\,d\theta$ near $p$ by \eqref{eq:volume element near p} and $O(t^5)\,dt\,d\theta$ near $q$ by \eqref{eq:volume element near q}. Since it is smooth on each compact subset of $M\setminus\{p,q\}$, it is therefore integrable on $\hm$, and $F_\infty<\infty$. The monotone convergence theorem implies $F(r)\rightarrow F_\infty$ as $r\rightarrow\infty$.

For sufficiently large $r$, Lemma \ref{lem:graphical expression} and \eqref{eq:volume element near q} give
\begin{align*}
F_\infty-F(r)&=\int_{\{\rho>r\}}V|\hn\rho|_{\hg}^2\dv\\
&=\int_{\S^3}\int_0^{\tau(r,\theta)}\left(\gamma^{-1}|W(q)|^2_g\,t^5+O(t^6)\right)\,dt\,d\theta.
\end{align*}
Noting $\tau(r,\theta)=\gamma^{-1/2}r^{-1}+O(r^{-2})$, we have
\begin{align*}
F_\infty-F(r)&=\int_{\S^3}\left(\frac{\gamma^{-1}|W(q)|_g^2}{6}\tau(r,\theta)^6+O(\tau(r,\theta)^7)\right)\,d\theta\\
&=\frac{\omega_3}{6}\gamma^{-1}|W(q)|^2_g\left(\gamma^{-3}r^{-6}+O(r^{-7})\right)+O(r^{-7})\\
&=\frac{\omega_3}{6}\gamma^{-4}|W(q)|^2_g r^{-6}+O(r^{-7}),
\end{align*}
which proves \eqref{eq:asymptotic of F at infinity}.
\end{proof}

We now prove Theorem \ref{thm1} using the ODE argument used in \cite{LP25}.

\begin{proof}[Proof of Theorem \ref{thm1}] We recall that
\begin{equation*}
|\hn\rho|_{\hg}\ge\delta\quad\text{on $M\setminus\{p,q\}$}\quad \text{and}\quad \sup_{\hm}\rho^2|\hw|_{\hg}\le c\delta^2
\end{equation*} by assumption. The coarea formula gives 
\begin{equation*}
F'(r)=\int_{\ls}V|\hn\rho|_{\hg}\da.
\end{equation*}
Using the divergence theorem and the first equation in \eqref{eq:lem2.1}, we have
\begin{equation*}
F'(r)=\int_{\sls}\left(\langle\hn V,\hn\rho\rangle+3\rho^{-1}V|\hn\rho|^2_{\hg}\right)\dv.
\end{equation*}
Differentiating by the coarea formula yields 
\begin{equation*}
F''(r)-\frac{3}{r}F'(r)=\int_{\ls}\partial_\nu V\da=\int_{\sls}\hd V\dv.
\end{equation*}
Next, integration by parts and the second equation in \eqref{eq:lem2.1} gives
\begin{align*}
\int_{\sls}\rho^2\hd V\dv&=r^2\int_{\ls}\partial_\nu V\da-\int_{\sls}\langle\hn\rho^2, \hn V\rangle\dv\\
&=r^2\left(F''-\frac{3}{r}F'\right)-2rF'+8F.
\end{align*}
Hence we have
\begin{equation*}
r^2F''(r)-5rF'(r)+8F(r)=\int_{\sls}\rho^2\hd V\dv.
\end{equation*}
Let 
\begin{equation*}
\omega=\sup_{\hm}\rho^2|\hw|_{\hg}.
\end{equation*}
Lemma \ref{lem:Kato}, together with the assumption \eqref{eq:assumption sec1}, yields that 
\begin{align*}
r^2F''-5rF'+8F&\ge-C\int_{\sls}\rho^2|\hw|_{\hg}V\dv\\&\ge -C\omega\int_{\sls}V\dv\\
&\ge-C\delta^{-2}\omega F(r).
\end{align*}
If we let $\Theta=C\delta^{-2}\omega$, then
\begin{equation}\label{eq:ODE}
r^2F''-5rF'+(8+\Theta)F\ge0.
\end{equation}

We now choose $c>0$ so that $0\le \Theta<1$; since $\Theta=C\delta^{-2}\omega$ and $\omega \le c\delta^2$ by hypothesis, it suffices to take $c=1/(2C)$, where $C$ is the universal constant of Lemma \ref{lem:Kato}. If $\Theta=0$, then $\omega=0$ and $\wh{W}\equiv0$, as desired. Thus, we assume $0<\Theta<1$. Let $x\ge y$ be the roots of the homogeneous indicial equation
\begin{equation*}
\lambda^2-6\lambda+(8+\Theta)=0,
\end{equation*}
that is,
\begin{equation*}
x=3+\sqrt{1-\Theta},\qquad y=3-\sqrt{1-\Theta},
\end{equation*}
and $2<y<3<x<4$. If we define
\begin{align*}
h(r)=\frac{F(r)}{r^x},
\end{align*}
then the inequality \eqref{eq:ODE} implies
\begin{equation}\label{eq:ODE2}
\left(\frac{h'}{r^{y-x-1}}\right)'=\left(\frac{rF'-xF}{r^y}\right)'=r^{-y-1}(r^2F''-5rF'+(8+\Theta)F)\ge0.
\end{equation}
By Proposition \ref{prop:asymptotic of F near 0}, we have
\begin{equation*}
\frac{rF'-xF}{r^y}=O(r^{4-y})\rightarrow0\quad\text{as } r\rightarrow0.
\end{equation*}
It then follows from \eqref{eq:ODE2} that $h'(r)\ge0$ for every $r>0$. On the other hand, Proposition \ref{prop:asymptotic of F at infinity} gives 
\begin{align*}
h(r)=\frac{F(r)}{r^x}\rightarrow0\quad\text{as } r\rightarrow\infty. 
\end{align*}
Since $h$ is non-negative and non-decreasing, it must hold that $h\equiv 0$, and hence $F\equiv 0$. By the lower bound \eqref{eq:assumption sec1} and the definition of $V$, we obtain $\hw\equiv0$ on $\hm$.

The conformal invariance of the Weyl tensor implies $W_g\equiv 0$. Since a locally conformally flat Einstein metric has constant sectional curvature, $(M,g)$ has constant sectional curvature 1. Then, by the Killing--Hopf theorem, $(M,g)$ is isometric to a spherical space form $\S^4/\Gamma$ for some finite group $\Gamma\subset \mathrm{O}(5)$ acting freely on $\S^4$.

In dimension four, the spherical space forms are $\S^4$ or $\mathbb{RP}^4$. To exclude the latter case, we claim that $|\hn\rho|_{\hg}=0$ at some points in the blow-up of $\mathbb{RP}^4$. Denote the covering map by $\pi:\S^4\rightarrow\mathbb{RP}^4$. Fix two distinct points $p,q\in\S^4$ and let $[p]=\pi(p),[q]=\pi(q)\in\RP^4$. Then we have
\begin{equation*}
\pi^\ast G_{[p]}^{\RP^4}=G_{p}^{\S^4}+G_{-p}^{\S^4}.
\end{equation*}
Since we know that $G^{\S^4}_p=(2\sin(d(p,\cdot)/2))^{-2}$ and $d(-p,\cdot)=\pi-d(p,\cdot)$, we have
\begin{equation*}
\pi^\ast G_{[p]}^{\RP^4}=\frac{1}{4\sin^2(d(p,\cdot)/2)}+\frac{1}{4\cos^2(d(p,\cdot)/2)}=\frac{1}{\sin^2d(p,\cdot)}.
\end{equation*}
Thus if we set $u=\langle x,p\rangle, v=\langle x,q\rangle$ on $\S^4\subset\R^5$, then $u^2,v^2$ have the same values at antipodal points, so are well-defined on $\RP^4$. Since we have $G_{[p]}^{\RP^4}=\frac{1}{1-u^2}$, it gives
\begin{equation*}
\rho^2=\frac{G_q}{G_p}=\frac{1-u^2}{1-v^2}.
\end{equation*}

Since $G_{[p]}^{\RP^4}=\frac{1}{1-u^2}$ has minimum on the cut locus $\mathrm{Cut}(p)=\{u=0\}\cong\RP^3$, where $d(p,\cdot)=\pi/2$, $\nabla G_{[p]}^{\RP^4}=0$ there. For the same reason, $\nabla G_{[q]}^{\RP^4}=0$ on $\mathrm{Cut}(q)=\{v=0\}$. Hence, we have $\nabla G_{[p]}^{\RP^4}=\nabla G_{[q]}^{\RP^4}=0$ on $\mathrm{Cut}(p)\cap\mathrm{Cut}(q)\cong\RP^2$, $\nabla\rho=0$ and $\hn\rho=0$ there. Therefore, we deduce $\Gamma$ is trivial and $(M,g)$ is isometric to $\S^4$.  
\end{proof}

\begin{remark}
Neither the maximality of $b_p$ at $q$ nor any relation between $p$ and $q$, except $p\neq q$, is used in this section. These enter only in the verification of the hypotheses in Theorem \ref{thm1} under the smallness of the integral gap to prove the gap theorem, where $q$ is chosen to be a maximum point of $b_p$, so that, in the limit, $p$ and $q$ become antipodal on the round sphere.
\end{remark}

\section{Integral Gap and Almost-Maximal Volume}
In this section we prove Theorem \ref{thm2} and deduce Corollaries \ref{cor3} and \ref{thm3}. Let $(M^n,g)$ be a closed connected Riemannian manifold of dimension $n\ge 3$ with $\ric\ge(n-1)g$. Fix $p\in M$, and let $G=G_p$ and $b=b_p$ be given by \eqref{eq:green-function-normalization}-\eqref{eq:function b}. Throughout this section, we write
\begin{equation*}
    s(r)=2\sin(r/2),\qquad c(r)=\cos(r/2),
\end{equation*}
so that the definition \eqref{eq:function b} reads
\begin{equation*}
G=s(b)^{2-n}.
\end{equation*}
Integrating the normalization \eqref{eq:green-function-normalization} over $M$ gives the identity
\begin{equation}\label{eq:integral of G}
\frac{n}{4}\int_M G\,dV_g=\omega_{n-1}.
\end{equation}

Set $m=\max_M b$. By \cite[Theorem 1.1]{Man25} we have $m\le \pi$, and if $m=\pi$, then $(M,g)$ is isometric to the round sphere, in which case Theorem \ref{thm2}, Corollaries \ref{cor3} and \ref{thm3} are trivial. We therefore assume that $m<\pi$ for the remainder of this section.

For a regular value $r\in(0,m)$ of $b$, we define 
\begin{equation}\label{eq:JSH}
J(r)=\frac{n}{4}\int_{\bsls}G\dvol,\quad S(r)=\int_{\bls}|\nabla b|\darea,\quad H(r)=\int_{\bls}\frac{1-|\nabla b|^2}{|\nabla b|}\darea,
\end{equation}
and extend the domain of definition of $S$ and $H$ to all of $(0,m)$ by assigning arbitrary values at critical values of $b$. By the gradient estimate \eqref{eq:manea}, $S,H\ge0$. Also, the function $J$ is nondecreasing with $J(0^+)=0$ and, by \eqref{eq:integral of G}, $J(m)=\omega_{n-1}$.

\subsection{Level set identities}
We first record that the critical set of $b$ is negligible.

\begin{lemma}\label{lem:critical set of b}
We have
\begin{equation*}
\mathcal H^n\left(\{x\in M\setminus\{p\}\,:\,|\nabla b|(x)=0\}\right)=0.
\end{equation*}
\end{lemma}

\begin{proof}
On every relatively compact coordinate ball in $M\setminus\{p\}$, the function $G$ is a nonconstant solution of the homogeneous uniformly elliptic equation $LG=0$. The volume estimate for critical sets of linear elliptic equations then implies that $\{|\nabla G|=0\}$ has zero $n$-dimensional measure \cite{NaberValtorta2017}. Since $m<\pi$, the function relating $b$ to $G$, given by $x \mapsto s(x)^{2-n}$, has nonvanishing derivative on $M\setminus\{p\}$; hence $b$ and $G$ have the same critical set there.
\end{proof}

\begin{lemma}\label{lem:coarea}
For every nonnegative Borel function $\varphi:(0,m)\rightarrow\R$,
\begin{equation}\label{eq:3.3}
\begin{aligned}
&\int_0^m \varphi(r)(S(r)+H(r))\,dr=\int_M\varphi(b)\,dV_g, \\
&\int_0^m\varphi(r) H(r)\,dr=\int_M\varphi(b)(1-|\nabla b|^2)\,dV_g.
\end{aligned}
\end{equation}
In particular, functions $J$ and $r\mapsto \vol(\{b\le r\})$ are absolutely continuous on $[0,m]$, and for a.e. $r\in(0,m)$,
\begin{equation}\label{eq:3.4}
J'(r)=\frac{n}{4}s(r)^{2-n}(S(r)+H(r)),\quad \frac{d}{dr}\vol(\bsls)=S(r)+H(r).
\end{equation}
\end{lemma}

\begin{proof}
The function $b$ is smooth on $M\setminus\{p\}$ with $0<b\le m$ there, and the level set $\{b=m\}$ consists of maximum points of $b$ and is hence contained in the critical set. Since the critical set has measure zero by Lemma \ref{lem:critical set of b}, the coarea formula gives, for every nonnegative Borel function $\psi$ on $M\setminus\{p\}$, 
\begin{equation}\label{eq:lem3.2-1}
\int_M\psi\dvol=\int_{\{|\nabla b|>0\}}\psi \dvol=\int_0^m\int_{\bls}\frac{\psi}{|\nabla b|}\darea\,dr.
\end{equation}
By Sard's theorem, a.e. $r\in(0,m)$ is a regular value of $b$, and at every regular value,
\begin{equation*}
\int_{\bls}\frac{\varphi(b)}{|\nabla b|}\darea=\varphi(r)\int_{\bls}\frac{|\nabla b|^2+(1-|\nabla b|^2)}{|\nabla b|}\darea=\varphi(r)(S(r)+H(r)).
\end{equation*}
Thus the choice $\psi=\varphi(b)$ in \eqref{eq:lem3.2-1} proves the first equation in \eqref{eq:3.3}.

Taking $\varphi=\frac{n}{4}s^{2-n}1_{(0,r]}$ and $\varphi=1_{(0,r]}$ in the first identity of \eqref{eq:3.3} and using $G=s(b)^{2-n}$, we obtain
\begin{equation*}
J(r)=\frac{n}{4}\int_0^r s^{2-n}(S+H)\,dt,\qquad \vol(\bsls)=\int_0^r(S+H)\,dt
\end{equation*}
for every $r\in(0,m)$. Since their values approach $\omega_{n-1}$ and $\vol(M)$ as $r \to m$, respectively, both integrands belong to $L^1(0,m)$, and the absolute continuity and \eqref{eq:3.4} therefore follow. 
\end{proof}

\begin{lemma}[cf. {\cite[(4.6)]{Man25}}]\label{eq:lemma 3.3}
For every regular value $r\in(0,m)$ of $b$,
\begin{equation}\label{eq:3.8}
S(r)=s(r)^{n-1}c(r)^{-1}(\omega_{n-1}-J(r)).
\end{equation}
\end{lemma}

\begin{proof}
Since $G=s(b)^{2-n}$ and $s'=c$, we have $\nabla G=(2-n)s(b)^{1-n}c(b)\nabla b$. Hence, on a regular level set $\bls$ with unit normal $\nu=\nabla b/|\nabla b|$, we have
\begin{equation}\label{eq:lem3.3-1}
\int_{\bls}\partial_\nu G\darea=(2-n)s(r)^{1-n}c(r)S(r).
\end{equation}
Since $\nabla G\neq 0$ near $p$ from the local expansion of the Green function at the pole, every sufficiently small $\epsilon>0$ is a regular value of $b$, and 
\begin{equation}\label{eq:lem3.3-2}
\int_{\{b=\epsilon\}}\partial_\nu G\darea\rightarrow -(n-2)\omega_{n-1}\quad\text{as $\epsilon\rightarrow0$}
\end{equation}
by the normalization \eqref{eq:green-function-normalization}. Since $\Delta G=\frac{n(n-2)}{4}G$ on $\{\epsilon\le b\le r\}$, the divergence theorem implies
\begin{equation*}
\frac{n(n-2)}{4}\int_{\{\epsilon\le b\le r\}}G \dvol=\int_{\bls}\partial_\nu G\darea-\int_{\{b=\epsilon\}}\partial_\nu G\darea.
\end{equation*}
Letting $\epsilon\rightarrow0$ along regular values and using \eqref{eq:lem3.3-1}, \eqref{eq:lem3.3-2}, and $J(0^+)=0$, we obtain
\begin{equation*}
(n-2)J(r)=(2-n)s(r)^{1-n}c(r)S(r)+(n-2)\omega_{n-1},
\end{equation*}
which is \eqref{eq:3.8}.
\end{proof}

\subsection{Almost-Maximal Volume and Gap Theorem}
The key quantity in the proof of Theorem \ref{thm2} is the following weighted integral gap:
\begin{equation}\label{eq:weighted integral gap}
\eta:=\frac{n}{4}\int_M Gc(b)^2(1-|\nabla b|^2)\dvol.
\end{equation}

\begin{proposition}\label{prop3.4}
There exists a constant $C_n>0$ such that
\begin{equation}\label{eq:prop3.4}
1-\frac{\vol(M,g)}{\vol(\S^n)}\le C_n\left(\frac{\eta}{\omega_{n-1}}\right)^{\frac{n}{n+2}}.
\end{equation}
\end{proposition}
\begin{proof}
We first note that by \eqref{eq:integral of G}
\begin{equation*}
\eta \le \frac{n}{4}\int_M G\dvol=\omega_{n-1}.
\end{equation*}
By Lemma \ref{lem:coarea}, 
\begin{equation}\label{eq:prop3.4-1}
\eta=\frac{n}{4}\int_0^mc(t)^2s(t)^{2-n}H(t)\,dt
\end{equation}
and
\begin{equation}\label{eq:prop3.4-2}
\frac{n}{4}\int_0^ms^{2-n}H\,dt=\frac{n}{4}\int_M G(1-|\nabla b|^2)\dvol\le\omega_{n-1}.
\end{equation}
Define 
\begin{equation*}
D(r)=J(r)-\omega_{n-1}(1-c(r)^n),\quad r\in[0,m].
\end{equation*}
Then, $D$ is absolutely continuous, $D(0)=0$ and $D(m)=\omega_{n-1}c(m)^n$ since $J(m)=\omega_{n-1}$. Substituting \eqref{eq:3.8} into \eqref{eq:3.4}, we find that, for a.e. $r\in(0,m)$,
\begin{equation*}
D'(r)+\frac{n}{2}\tan(r/2)D(r)=\frac{n}{4}s(r)^{2-n}H(r)
\end{equation*}
and
\begin{equation*}
D(r)=\frac{n}{4}c(r)^n\int_0^rc(t)^{-n}s(t)^{2-n}H(t)\,dt\ge0.
\end{equation*}
Since $c$ is decreasing,
\begin{equation*}
c(t)^{-n}=c(t)^2c(t)^{-(n+2)}\le c(t)^2c(r)^{-(n+2)}
\end{equation*}
for $0<t\le r$, and hence, by \eqref{eq:prop3.4-1} and \eqref{eq:prop3.4-2}, 
\begin{equation}\label{eq:3.13}
0\le D(r)\le \eta c(r)^{-2},\quad r\in(0,m].
\end{equation}
Taking $r=m$ in \eqref{eq:3.13} gives
\begin{equation}\label{eq:3.14}
\omega_{n-1}c(m)^{n+2}\le \eta.
\end{equation}
In particular, $\eta>0$, since $m<\pi$ and hence $c(m)>0$.

Next, by Lemma \ref{lem:coarea} and Lemma \ref{eq:lemma 3.3}, we have, for every $r\in(0,m)$,
\begin{align*}
\vol(M)&\ge\vol(\bsls)=\int_0^r(S+H)\,dt\\
&\ge\int_0^rs(t)^{n-1}c(t)^{-1}(\omega_{n-1}-J(t))\,dt\\
&=\int_0^rs(t)^{n-1}c(t)^{-1}(\omega_{n-1}c(t)^n-D(t))\,dt\\
&=\underbrace{\omega_{n-1}\int_0^rs(t)^{n-1}c(t)^{n-1}\,dt}_{:=A}-\underbrace{\int_0^rs(t)^{n-1}c(t)^{-1}D(t)\,dt}_{:=B}.
\end{align*}
The first term $A$ can be computed explicitly:
\begin{equation}\label{eq3.15}
A=\omega_{n-1}\int_0^rs(t)^{n-1}c(t)^{n-1}\,dt=\omega_{n-1}\int_0^r\sin^{n-1}t\,dt=\vol\left(B_r^{\S^n}\right),
\end{equation}
where $\vol\left(B_r^{\S^n}\right)$ denotes the volume of the geodesic ball of radius $r$ in $\S^n$. The second term $B$ can be controlled by \eqref{eq:3.13}: 
\begin{equation}\label{eq3.16}
\begin{aligned}
B&=\int_0^rs(t)^{n-1}c(t)^{-1}D(t)\,dt\le \eta\int_0^rs(t)^{n-1}c(t)^{-3}\,dt\\
&\le 2^{n-2}\eta\int_0^rs(t)c(t)^{-3}\,dt=2^{n-1}\eta (c(r)^{-2}-1)\le 2^{n-1}\eta c(r)^{-2}.
\end{aligned}
\end{equation}
On the other hand, we note 
\begin{equation}\label{eq3.17}
\vol\left(\S^n\setminus B_r^{\S^n}\right)=\omega_{n-1}\int_r^{\pi}\sin^{n-1}t\,dt\le\frac{2^{n}\omega_{n-1}}{n}c(r)^n.
\end{equation}
Combining \eqref{eq3.15}, \eqref{eq3.16} and \eqref{eq3.17} gives 
\begin{equation}\label{eq3.18}
\vol(\S^n)-\vol(M)\le \frac{2^n\omega_{n-1}}{n}c(r)^n+2^{n-1}\eta c(r)^{-2},\quad r\in(0,m].
\end{equation}

Finally, we set
\begin{equation*}
    \beta=\left(\frac{\eta}{\omega_{n-1}}\right)^{\frac{1}{n+2}}\in(0,1].
\end{equation*}
If $\beta\ge\frac{1}{2}$, then \eqref{eq:prop3.4} holds with $C_n=2^n$ since its left-hand side is at most $1\le 2^n\beta^n$. If $\beta<\frac{1}{2}$, then $c(m)\le \beta<1$ by \eqref{eq:3.14}, and by continuity, there exists $r_\ast\in(0,m]$ with $c(r_\ast)=\beta$. Then \eqref{eq3.18} at $r=r_\ast$ gives 
\begin{equation*}
\vol(\S^n)-\vol(M)\le\left(\frac{2^n}{n}+2^{n-1}\right)\omega_{n-1}\left(\frac{\eta}{\omega_{n-1}}\right)^{\frac{n}{n+2}}
\end{equation*}
since $\eta\beta^{-2}=\omega_{n-1}\beta^n$. It completes the proof.
\end{proof}

We now prove Theorem \ref{thm2}.

\begin{proof}[Proof of Theorem \ref{thm2}]
Let $h$ be the unique solution to
\begin{equation}\label{eq3.19}
Lh=nc(b)^2(1-|\nabla b|^2)\quad\text{on $M$}.
\end{equation}
Note that $h\in C^1(M)$, and $h\ge0$ if we test \eqref{eq3.19} with $\min\{h,0\}$.

By the normalization \eqref{eq:green-function-normalization} and the symmetry of the Green function,
\begin{equation}\label{eq3.20}
h(p)=\frac{1}{(n-2)\omega_{n-1}}\int_M Gnc(b)^2(1-|\nabla b|^2)\dvol=\frac{4\eta}{(n-2)\omega_{n-1}}.
\end{equation}
Integrating \eqref{eq3.19} over $M$, we have
\begin{equation}\label{eq3.21}
\begin{aligned}
\fint_Mh\dvol&=\frac{4}{n(n-2)}\fint_M nc(b)^2(1-|\nabla b|^2)\dvol\\
&\le\frac{4}{n-2}\fint_M (1-|\nabla b|^2)\dvol=\frac{4}{n-2}\E_p(M,g).
\end{aligned}
\end{equation}
Since $h\ge0 $, we obtain the pointwise bound, together with \eqref{eq3.19},
\begin{equation*}
\Delta h=\frac{n(n-2)}{4}h-nc(b)^2(1-|\nabla b|^2)\ge -nc(b)^2(1-|\nabla b|^2).
\end{equation*}
The Neumann-type maximum principle due to Wei--Ye \cite[Theorem C]{WeiYe2007} gives a constant $C_{n,q}>0$ for each $q>n/2$ such that
\begin{equation*}
\sup_Mh\le\fint_Mh\dvol+C_{n,q}\left(\fint_M\left(nc(b)^2(1-|\nabla b|^2)\right)^q\right)^{1/q}.
\end{equation*}
Since $0\le c(b)^2(1-|\nabla b|^2)\le 1$, for each $q\ge 1$ the integrand of the last term is bounded by $n^q(1-|\nabla b|^2)$. Consequently the second term is bounded by $C_{n,q}n\E_p(M,g)^{1/q}$. Combining this with \eqref{eq3.20} and \eqref{eq3.21}, we obtain
\begin{equation}\label{eq3.22}
\begin{aligned}
\eta&=\frac{(n-2)\omega_{n-1}}{4}h(p)\le\frac{(n-2)\omega_{n-1}}{4}\sup_M h\\
&\le \omega_{n-1}\left(\E_p(M,g)+\frac{n(n-2)}{4}C_{n,q}\E_p(M,g)^{1/q}\right)\\
&\le C_{n,q}'\omega_{n-1}\E_p(M,g)^{1/q}
\end{aligned}
\end{equation}
since $0\le\E_p(M,g)\le 1$, where $C_{n,q}'=1+n(n-2)C_{n,q}/4$.

Finally, for given $\gamma<\frac{2}{n+2}$, if we set $q=\frac{n}{\gamma(n+2)}$ so that $q>\frac{n}{2}$, then Proposition \ref{prop3.4} and \eqref{eq3.22} give
\begin{equation*}
1-\frac{\vol(M,g)}{\vol(\S^n)}\le C_n\left(\frac{\eta}{\omega_{n-1}}\right)^{\frac{n}{n+2}}\le C_{n,\gamma}\E_p(M,g)^{\frac{n}{q(n+2)}}=C_{n,\gamma}\E_p(M,g)^{\gamma},
\end{equation*}
as desired. This completes the proof.
\end{proof}

We now prove the two corollaries stated in the introduction.

\begin{proof}[Proof of Corollary \ref{cor3}]
Fix $0<\gamma<\frac{2}{n+2}$. By Theorem \ref{thm2}, 
\begin{equation*}
\vol(\S^n)-\vol(M,g)\le C_{n,\gamma}\vol(\S^n)\E_p(M,g)^{\gamma}.
\end{equation*}
Hence, if $\ve(n)$ is sufficiently small, the volume of $(M,g)$ is almost maximal, and Aubry's volume pinching theorem \cite{Aubry2005}, applied with $k=1$, shows that $\pi_1(M)$ is trivial, that $M$ is diffeomorphic to $\S^n$, and that
\begin{equation*}
d_{\mathrm{GH}}((M,g),(\S^n,g_{\S^n}))\le C(n)(\vol(\S^n)-\vol(M,g))^{\beta(n)},
\end{equation*}
where $C(n)>0$ and $\beta(n)>0$ are the constants of \cite{Aubry2005}. Combining the two estimates proves the corollary with $\alpha_n=\gamma\beta(n)$. 
\end{proof}

\begin{proof}[Proof of Corollary \ref{thm3}]
By Theorem \ref{thm2}, if $\E_p(M,g)\le \ve(n)$ with $\ve(n)$ sufficiently small, then $\vol(M,g)$ is arbitrarily close to $\vol(\S^n)$. Since we assumed that $(M,g)$ is Einstein with $\ric=(n-1)g$, the sphere theorem of Honda--Mondello \cite[Theorem C]{HondaMondello} implies that $(M,g)$ is isometric to the round sphere $(\S^n,g_{\S^n})$. 
\end{proof}

\begin{remark}\label{rmk3.5}
By \cite[Proposition 6.1]{Man25} with notations used there, 
\begin{equation*}
\eta=\omega_{n-1}-A(m)=\omega_{n-1}-\frac{n(n+2)}{32}\int_{M}G^{\frac{n-4}{n-2}}\dvol.
\end{equation*}
When $n=4$, the exponent $(n-4)/(n-2)$ vanishes so that
\begin{equation}\label{eq:eta}
\eta=\omega_3\left(1-\frac{\vol(M,g)}{\vol(\S^4)}\right).
\end{equation}
Thus, Proposition \ref{prop3.4} holds with exponent 1 in place of $\frac{n}{n+2}$ when $n=4$, and Theorem \ref{thm2} holds with $\gamma<\frac{1}{2}$.

The identity \eqref{eq:eta} also agrees with the formulation of Gursky--Malchiodi \cite{GurskyMalchiodi}. When $n=4$, let $\wh{g}=G^2g$ be the conformal blow-up on $M\setminus\{p\}$ (with pole at $p$) and let $F=G^{-1}(1+|\hn G|_{\hg}^2)-4$ be the function introduced in \cite[(3.1)]{GurskyMalchiodi}. Substituting $G=s(b)^{-2}$ gives the pointwise identity $-F=4c(b)^2(1-|\nabla b|^2)$. Hence the estimate \eqref{eq:manea} is equivalent to the estimate $|\hn G|_{\hg}^2\le 4G-1$ of \cite[Proposition 3.3]{GurskyMalchiodi} in dimension four, and the weighted integral gap $\eta$ becomes 
\begin{equation*}
\eta=-\frac{1}{4}\int_M FG\dvol.
\end{equation*}
So the weight $c(b)^2$ in \eqref{eq:weighted integral gap} is natural in light of \cite{GurskyMalchiodi}. Moreover, for Einstein metrics, combining the mass formula \cite[Theorem 1.1]{GurskyMalchiodi} with the identity above, we obtain that $\eta\le 2\pi^2 m(\hg)$, where $m(\hg)$ denotes the ADM mass of the blow-up.
\end{remark}

\section{Another Proof of Four-Dimensional Gap Theorem}
In this section we prove Theorem \ref{thm4}. We verify that the two hypotheses of Theorem \ref{thm1} are satisfied when the normalized integral gap is small, thereby providing another proof of Corollary \ref{thm3} in dimension four without using the sphere theorem of Honda--Mondello \cite{HondaMondello}.

\begin{proof}[Proof of Theorem \ref{thm4}]
To the contrary, suppose that the conclusion fails for some fixed $0<\delta<2$ and $\zeta>0$. Then there exists a sequence $(M_i,g_i)$ of closed connected Einstein 4-manifolds with $\ric_{g_i}=3g_i$, together with points $p_i,q_i\in M_i$, such that
\begin{equation*}
q_i\in\underset{x\in M_i}{\mathrm{argmax}}\,b_{p_i}(x)\quad\text{and}\quad \E_{p_i}(M_i,g_i)\rightarrow0,
\end{equation*}
but, for every $i$, at least one of the following inequalities fails: 
\begin{equation}\label{eq:thm1.4.1}
\inf_{\hm_i\setminus\{p_i\}}|\hn\rho_{i}|_{\hg_i}\ge\delta,\qquad \sup_{\hm_i}\rho_i^2|\hw_i|_{\hg_i}<\zeta,
\end{equation}
where
\begin{equation*}
\hg_i:=G_{q_i}^2 g_i\quad\text{on $M_i\setminus\{q_i\}$}\quad\text{and}\quad \rho_i:=\left(\frac{G_{q_i}}{G_{p_i}}\right)^{1/2}.
\end{equation*}

By Theorem \ref{thm2} and Bishop--Gromov comparison,
\begin{equation*}
\vol(M_i,g_i)\rightarrow\vol(\S^4).
\end{equation*}
Then Colding's almost-maximal-volume theorem \cite{Colding1996Shape} implies that $(M_i,g_i)$ converges to $\S^4$ in the Gromov-Hausdorff sense. Since $(M_i,g_i)$ are Einstein, up to taking a further subsequence (which we again index by $i$), there exist diffeomorphisms $\Phi_i:\S^4\rightarrow M_i$ such that $\Phi_i^\ast g_i\rightarrow g_{\S^4}$ smoothly, $\Phi_i^{-1}(p_i)\rightarrow p_\infty$, and $\Phi_i^{-1}(q_i)\rightarrow q_\infty$; see for instance \cite{Anderson1990}. The Green functions also converge smoothly on compact sets away from their poles, while their pole expansions are uniform under the convergence; see \cite{LeeParker1987}. If $p_\infty^\ast$ denotes the antipodal point of $p_\infty$, then we have
\begin{equation*}
b_{p_i}(\Phi_i(p_\infty^\ast))\rightarrow\pi.
\end{equation*}
Since $b_{p_i}\le\pi$ and $q_i$ is a maximum point of $b_{p_i}$, we conclude that $b_{p_i}(q_i)\rightarrow\pi$. Since $b_{p_i}$ is 1-Lipschitz and $b_{p_i}(p_i)=0$, we have
\begin{equation*}
b_{p_i}(q_i)\le d_{g_i}(p_i,q_i)\le\pi
\end{equation*}
and $d_{g_i}(p_i,q_i)\rightarrow\pi$. Therefore $q_\infty=p_\infty^\ast$. 

First we consider the Weyl term. The conformal transformation law and the definition of $\rho_i$ yield
\begin{equation*}
\rho_i^2|\wh{W}_i|_{\wh{g}_i}=\frac{|W_i|_{g_i}}{G_{p_i}G_{q_i}}.
\end{equation*}
Note that in dimension four, we have $G_a=(2\sin(b_a/2))^{-2}\ge1/4$ for each $a\in M$. Thus
\begin{equation}
\sup_{\hm_i}\rho_i^2|\hw_i|_{\hg_i}\le16\norm{W_{g_i}}_{L^\infty(M_i,g_i)}\rightarrow0,
\end{equation}
since the limit metric is round and hence has vanishing Weyl tensor. In particular, the second inequality in \eqref{eq:thm1.4.1} holds for all sufficiently large $i$.

It remains to prove a uniform lower bound for $|\hn\rho_i|_{\hg_i}$. On the limiting sphere, denoting $r=d_{\S^4}(p_\infty,\cdot)$, we have
\begin{equation*}
G_{p_\infty}=(2\sin(r/2))^{-2},\quad G_{q_\infty}=(2\cos(r/2))^{-2},\quad \rho_\infty=\tan(r/2)
\end{equation*}
and so
\begin{equation}\label{eq:thm1.4.2}
|\hn\rho_\infty|_{\hg_\infty}=G_{q_\infty}^{-1}|\nabla\rho_\infty|_{g_{\S^4}}=2 \quad\text{on $\S^4\setminus\{p_\infty,q_\infty\}$}.
\end{equation}

Set 
\begin{equation*}
    \gamma_i:=G_{q_i}(p_i)=G_{p_i}(q_i).
\end{equation*}
By smooth convergence of Green functions and noting $q_\infty=p_\infty^\ast$, we have
\begin{equation}\label{eq:thm1.4.3}
\gamma_i\rightarrow G_{q_\infty}(p_\infty)=\frac{1}{4}.
\end{equation}
Let $s_i=d_{g_i}(p_i,\cdot)$. The uniform pole expansions give, near $p_i$,
\begin{equation*}
G_{p_i}=s_i^{-2}(1+O_1(s_i^2)),\qquad G_{q_i}=\gamma_i+O_1(s_i),
\end{equation*}
where the constants in the $O_1$ terms are independent of $i$ for all sufficiently large $i$. Hence
\begin{equation*}
\rho_i=\gamma_i^{1/2}s_i+O_1(s_i^2),\qquad G_{q_i}^{-1}=\gamma_i^{-1}+O(s_i),
\end{equation*}
and 
\begin{equation}\label{eq:thm1.4.4}
|\hn\rho_i|_{\hg_i}=G_{q_i}^{-1}|\nabla\rho_i|_{g_i}=\gamma_i^{-1/2}+O(s_i)
\end{equation}
near $p_i$, uniformly in $i$.

Similarly, if $t_i=d_{g_i}(q_i,\cdot)$, then near $q_i$,
\begin{equation*}
G_{q_i}=t_i^{-2}(1+O_1(t_i^2)),\qquad G_{p_i}=\gamma_i+O_1(t_i),
\end{equation*}
so that
\begin{equation*}
\rho_i=\gamma_i^{-1/2}t_i^{-1}+O_1(1),\qquad G_{q_i}^{-1}=t_i^2(1+O(t_i^2)).
\end{equation*}
It then follows that
\begin{equation}\label{eq:thm1.4.5}
|\hn\rho_i|_{\hg_i}=\gamma_i^{-1/2}+O(t_i)
\end{equation}
near $q_i$, again uniformly in $i$.

By \eqref{eq:thm1.4.3}, \eqref{eq:thm1.4.4} and \eqref{eq:thm1.4.5}, for $r_0>0$ sufficiently small and $i$ sufficiently large, we have
\begin{equation}\label{eq:thm1.4.7}
|\hn\rho_i|_{\hg_i}>\delta\quad\text{on } (B_{r_0}^{g_i}(p_i)\setminus\{p_i\})\cup(B_{r_0}^{g_i}(q_i)\setminus\{q_i\}). 
\end{equation}
On the compact complement of these two balls, the Green functions and their first derivatives converge smoothly to the Green function on $\S^4$. By \eqref{eq:thm1.4.2},
\begin{equation}\label{eq:thm1.4.8}
|\hn\rho_i|_{\hg_i}\rightarrow 2\quad\text{uniformly on } M_i\setminus(B_{r_0}^{g_i}(p_i)\cup B_{r_0}^{g_i}(q_i)).
\end{equation}
Since $\delta<2$, \eqref{eq:thm1.4.7} and \eqref{eq:thm1.4.8} show that the first inequality in \eqref{eq:thm1.4.1} also holds for all sufficiently large $i$, which is a contradiction. This proves Theorem \ref{thm4}. 
\end{proof}

Finally, we remark that taking any $0<\delta<2$ and then choosing $\zeta<c\delta^2$, Theorems \ref{thm1} and \ref{thm4} together yield another proof of Corollary \ref{thm3} in dimension 4.

\bibliographystyle{alpha}
\bibliography{reference2}
\end{document}
